\chapter{Metodologia}\label{cap_3}

Este capítulo descreve o funcionamento do \ac{FLIM}, os três componentes avaliados sobre ele, e o protocolo experimental. A descrição é detalhada de propósito: boa parte das conclusões desta dissertação decorre de propriedades específicas do procedimento de estimação dos filtros, e elas só ficam visíveis quando o procedimento é apresentado com precisão.

\section{O FLIM em detalhe}

\subsection{O que é um filtro convolucional}

Uma rede convolucional reconhece padrões por meio de filtros. Um filtro é uma pequena matriz de números, tipicamente de tamanho $3 \times 3$, que é deslizada sobre a imagem. Em cada posição, calcula-se a soma dos produtos entre os valores do filtro e os valores da imagem naquela vizinhança. O resultado é alto onde a vizinhança se parece com o padrão guardado no filtro, e baixo onde não se parece.

Formalmente, para uma imagem $I$ e um filtro $K$ de tamanho $h \times w$ aplicado na posição $(x, y)$:

\begin{equation}
(I * K)(x, y) = \sum_{i=1}^{h} \sum_{j=1}^{w} I(x + i - c_h,\; y + j - c_w) \cdot K(i, j)
\label{eq_conv}
\end{equation}

onde $c_h$ e $c_w$ indicam o centro do filtro. O conjunto de filtros de uma camada define o vocabulário visual com que a rede descreve a imagem, e uma rede com várias camadas compõe padrões simples das primeiras camadas em padrões mais complexos nas seguintes.

A diferença entre os métodos está em como esses filtros são obtidos. Redes treinadas por retropropagação os inicializam aleatoriamente e os ajustam ao longo de milhares de iterações, calculando o quanto cada valor contribuiu para o erro e corrigindo-o na direção oposta. Esse procedimento requer grandes conjuntos anotados.

\subsection{Estimação dos filtros a partir de marcadores}

O \ac{FLIM} não ajusta filtros, ele os \emph{extrai} \cite{souza2020flim}. O especialista desenha marcadores sobre poucas imagens, indicando regiões de objeto e de fundo. O método recorta as vizinhanças centradas nos pixels marcados e usa essas vizinhanças, após agrupamento, como filtros.

Seja $M$ o conjunto de pixels marcados em uma imagem. Para cada pixel $p \in M$, define-se o \emph{patch} $\phi(p)$ como o vetor obtido empilhando os valores das bandas na vizinhança $h \times w$ de $p$:

\begin{equation}
\phi(p) \in \mathbb{R}^{h \cdot w \cdot b}
\label{eq_patch}
\end{equation}

onde $b$ é o número de bandas da entrada daquela camada. Para a primeira camada, $b$ é o número de canais da imagem; para as seguintes, é o número de filtros da camada anterior.

O procedimento ocorre em dois níveis de agrupamento.

\paragraph{Primeiro nível, por marcador.} Os marcadores são separados em componentes conexas: um traço isolado do pincel é uma componente, e dois traços que se encostam formam uma só. Para cada componente $m$, os patches dos seus pixels são agrupados por $k$-médias em $n_m$ grupos, e os centróides resultantes tornam-se candidatos a filtro:

\begin{equation}
C_m = \operatorname{kmeans}\big(\{\phi(p) : p \in m\},\; n_m\big)
\label{eq_nivel1}
\end{equation}

O parâmetro $n_m$, chamado \texttt{nkernels\_per\_marker} na implementação, é definido na arquitetura. Nos experimentos deste trabalho ele vale 3 no conjunto de parasitas, 5 no BraTS e 15 no conjunto de conjuntivite.

\paragraph{Segundo nível, global.} Os candidatos de todas as componentes e de todas as imagens anotadas são reunidos em um único conjunto $C = \bigcup_m C_m$, que então é reduzido ao número de canais de saída da camada, $n_{\text{out}}$:

\begin{equation}
K = \operatorname{kmeans}(C,\; n_{\text{out}})
\label{eq_nivel2}
\end{equation}

Os centróides resultantes, após normalização, são os filtros da camada. A saída da camada serve de entrada para o mesmo procedimento na camada seguinte, de modo que a estimação avança camada a camada sem nenhum passo de retropropagação.

\subsection{Os dois regimes do segundo agrupamento}

A Equação \ref{eq_nivel2} tem um comportamento que é central para os resultados deste trabalho. A implementação só executa o agrupamento quando há mais candidatos do que vagas:

\begin{equation}
K =
\begin{cases}
C, & \text{se } |C| \leq n_{\text{out}} \\[4pt]
\operatorname{kmeans}(C, n_{\text{out}}), & \text{se } |C| > n_{\text{out}}
\end{cases}
\label{eq_regimes}
\end{equation}

Quando o número de candidatos não excede o número de vagas, os próprios candidatos tornam-se os filtros, sem agrupamento algum, e o número de canais da camada é reduzido para $|C|$. Acrescentar um marcador, nesse regime, acrescenta filtros e deixa os anteriores intactos.

Quando o número de candidatos excede as vagas, o agrupamento precisa escolher como particionar o conjunto. Acrescentar um marcador acrescenta candidatos, a partição é refeita sobre o conjunto ampliado, e os centróides se deslocam. Nenhum filtro é preservado por construção.

A razão entre candidatos e vagas difere bastante entre os conjuntos de dados usados aqui. Com cerca de cinquenta componentes de marcador distribuídas em três imagens, o conjunto de parasitas produz aproximadamente 150 candidatos para 200 vagas, operando abaixo do limite. O BraTS produz cerca de 250 candidatos para 8 vagas, e o conjunto de conjuntivite cerca de 750 para 16. Nos dois últimos casos há dezenas de candidatos disputando cada vaga, e o agrupamento é intenso.

\subsection{O decodificador adaptativo}\label{sec_decoder}

O codificador produz mapas de ativação, não uma máscara. A conversão é feita por um decodificador adaptativo \cite{soares2026flim}, que atribui a cada canal de ativação um peso em $\{-1, +1\}$ conforme uma heurística que estima se aquele canal responde ao objeto ou ao fundo, e combina os canais por uma convolução ponto a ponto:

\begin{equation}
S(x, y) = \sum_{c=1}^{n_{\text{out}}} w_c \cdot A_c(x, y), \qquad w_c \in \{-1, +1\}
\label{eq_decoder}
\end{equation}

onde $A_c$ é o mapa de ativação do canal $c$. Os pesos são recalculados para cada imagem de entrada, o que dá ao decodificador o caráter adaptativo. Nenhum parâmetro é ajustado por gradiente.

O mapa $S$ é normalizado para o intervalo $[0, 1]$ e binarizado por limiarização de Otsu, seguida de um filtro que remove componentes conexas fora de uma faixa de área esperada para o objeto.

\section{Métricas}

\subsection{Qualidade da segmentação}

A métrica principal é o \fbeta, que combina precisão e revocação com peso maior para a precisão, seguindo o trabalho de referência:

\begin{equation}
F_\beta = \frac{(1 + \beta^2) \cdot P \cdot R}{\beta^2 \cdot P + R}, \qquad \beta^2 = 0{,}3
\label{eq_fbeta}
\end{equation}

onde $P$ é a precisão e $R$ a revocação, calculadas pixel a pixel sobre a máscara binarizada.

Reporta-se também a \ac{IoU}, razão entre a interseção e a união das máscaras predita e verdadeira:

\begin{equation}
\text{IoU} = \frac{|\hat{Y} \cap Y|}{|\hat{Y} \cup Y|}
\label{eq_iou}
\end{equation}

A acurácia pixel a pixel é reportada por completude, mas não discrimina neste domínio: como os objetos ocupam pequena fração da imagem, predizer tudo como fundo já produz acurácia superior a $0{,}97$.

\subsection{Esforço de interação}

Para a avaliação interativa usa-se o \ac{NoC}, número de cliques necessários para que a \ac{IoU} ultrapasse um alvo $\tau$, com um teto $N_{\max}$ de interações:

\begin{equation}
\text{NoC@}\tau = \min\{\,n \leq N_{\max} : \text{IoU}_n \geq \tau\,\},
\quad \text{ou } N_{\max} \text{ se nenhum } n \text{ satisfaz}
\label{eq_noc}
\end{equation}

Imagens que não atingem o alvo entram na média com o valor do teto. Por isso a taxa de sucesso é reportada ao lado do \ac{NoC} em todas as tabelas: um método poderia exibir contagem baixa simplesmente por desistir antes, e a contagem isolada premiaria esse comportamento.

\section{Aprendizado ativo}

Os critérios de aquisição avaliados pertencem a duas famílias. Seja $P(x,y) \in (0,1)$ a probabilidade predita de o pixel $(x,y)$ pertencer ao objeto.

\subsection{Critérios de incerteza}

A \textbf{entropia} mede a indecisão do modelo, atingindo o máximo quando a probabilidade predita é $0{,}5$:

\begin{equation}
H(x,y) = -\big[P \log P + (1-P)\log(1-P)\big]
\label{eq_entropia}
\end{equation}

O \textbf{least confidence} mede a distância à decisão, com a mesma intuição e escala diferente:

\begin{equation}
\text{LC}(x,y) = 1 - |2P - 1|
\label{eq_lc}
\end{equation}

Para pontuar uma \emph{região} $r$, toma-se a média do critério sobre seus pixels:

\begin{equation}
s(r) = \frac{1}{|r|} \sum_{(x,y) \in r} H(x,y)
\label{eq_score_regiao}
\end{equation}

\subsection{Critérios de diversidade}

Os critérios de diversidade não olham a incerteza, e sim a representatividade no espaço de atributos. Cada candidato é descrito por um vetor obtido pela média das ativações do codificador dentro da sua máscara.

O \textbf{CoreSet} seleciona o candidato mais distante do que já foi anotado, que é o passo guloso do $k$-center:

\begin{equation}
r^{\ast} = \arg\max_{r \notin A} \; \min_{a \in A} \; \lVert f(r) - f(a) \rVert_2
\label{eq_coreset}
\end{equation}

onde $A$ é o conjunto já anotado e $f(\cdot)$ o vetor de atributos. O \textbf{medoide} faz o oposto, escolhendo o candidato mais típico:

\begin{equation}
r^{\ast} = \arg\min_{r} \; \frac{1}{|R|}\sum_{q \in R} \lVert f(r) - f(q) \rVert_2
\label{eq_medoide}
\end{equation}

Esses dois critérios são de imagem por construção na literatura, pois comparam imagens inteiras. A extensão ao nível de região, feita neste trabalho, substitui a média global das ativações pela média dentro da máscara de cada região, mantendo o mesmo espaço.

\subsection{Geração dos candidatos}

A pontuação só ordena a lista de candidatos que recebe, de modo que o procedimento que gera essa lista é parte do método, e não detalhe de implementação. Dois geradores são comparados nesta dissertação.

O primeiro, usual na literatura, particiona a imagem em superpixels pelo algoritmo \ac{SLIC}, e cada superpixel é um candidato. O segundo, proposto aqui, define cada candidato como um disco de raio $\rho$ centrado em um ponto do contorno predito:

\begin{equation}
\partial \hat{Y} = \hat{Y} \setminus \operatorname{erosão}(\hat{Y}), \qquad
r_i = \{(x,y) : \lVert (x,y) - c_i \rVert_2 \leq \rho,\; c_i \in \partial \hat{Y}\}
\label{eq_contorno}
\end{equation}

onde $\hat{Y}$ é a máscara predita. Os centros $c_i$ são amostrados sobre o contorno com espaçamento mínimo de $2\rho$, para que os discos se toquem sem se sobrepor. O contorno usado é o \emph{predito}, e não o verdadeiro, de modo que o procedimento não depende de anotação prévia.

\section{Aprendizado contínuo}

\subsection{O esquecimento no FLIM}\label{sec_cl_flim}

Os métodos consolidados de \ac{CL} supõem um modelo treinado por gradiente, e atuam sobre o deslocamento dos pesos durante o ajuste. Como o \ac{FLIM} não possui otimizador, nenhum deles se aplica diretamente.

O fenômeno análogo, contudo, está na Equação \ref{eq_nivel2}. Sejam $C_t$ o conjunto de candidatos na rodada $t$ e $K_t$ o banco de filtros resultante. Acrescentar marcadores produz $C_{t+1} \supset C_t$, e como o agrupamento é refeito sobre o conjunto ampliado, não há garantia de que $K_{t+1}$ preserve qualquer elemento de $K_t$. Além disso, $K_t$ é compartilhado por todas as imagens, inclusive as não anotadas, de modo que a perturbação é global.

\subsection{O mecanismo de retenção}

A intervenção proposta interpola o banco em vez de substituí-lo. Para cada centróide novo $k \in K_{t+1}$, localiza-se o centróide anterior mais próximo e produz-se uma combinação convexa:

\begin{equation}
\tilde{k} = (1 - \alpha)\,\operatorname{anc}(k) + \alpha\,k,
\qquad
\operatorname{anc}(k) = \arg\min_{k' \in K_t} \lVert k - k' \rVert_2
\label{eq_retencao}
\end{equation}

O parâmetro $\alpha \in [0,1]$ é a taxa de plasticidade. Com $\alpha = 1$ o procedimento reproduz o \ac{FLIM} original. Com $\alpha = 0$ o banco passa a conter apenas direções já presentes em $K_t$, e nenhuma informação nova entra. Valores intermediários deslocam o banco parcialmente na direção da nova estimativa.

A âncora é tomada por centróide \emph{novo}, e não por centróide anterior. A direção importa: ancorar no sentido oposto obrigaria a saída a ter o tamanho do banco anterior, o que conflita com o ajuste automático do número de canais descrito na Equação \ref{eq_regimes} e invalida as estatísticas de normalização recalculadas a cada rodada.

\section{Aprendizado interativo}

\subsection{O usuário simulado}

Avaliar refinamento interativo exige um procedimento reprodutível que decida onde o especialista clicaria. Adota-se o protocolo consolidado por \cite{xu2016deep_arxiv}, no qual o clique vai ao centro do maior erro.

O erro é separado em dois tipos. O falso negativo, $\text{FN} = Y \setminus \hat{Y}$, indica região de objeto que o modelo não detectou e gera clique de objeto. O falso positivo, $\text{FP} = \hat{Y} \setminus Y$, indica região de fundo marcada como objeto e gera clique de fundo. Entre as componentes conexas desses dois conjuntos, escolhe-se a de maior área.

Dentro da componente escolhida, o clique \textbf{não} vai ao centróide. Vai ao ponto de máxima transformada de distância:

\begin{equation}
c^{\ast} = \arg\max_{(x,y) \in E} \; D_E(x,y)
\label{eq_clique}
\end{equation}

onde $E$ é a componente de erro e $D_E(x,y)$ é a distância euclidiana de $(x,y)$ à borda de $E$. Essa escolha não é arbitrária. O centróide de uma região côncava pode cair fora dela, o que produziria um clique com rótulo errado; e mesmo em região convexa o centróide não é necessariamente o ponto mais afastado da borda. O ponto de máxima distância está, por construção, no interior e o mais longe possível da franja, onde a atribuição de rótulo seria ambígua.

O clique é materializado como um disco de raio 5 pixels, valor medido nos 31 conjuntos de marcadores reais disponíveis, recortado de modo a não ultrapassar a fronteira do objeto. Não há sorteio em nenhuma etapa do laço, de forma que o \ac{NoC} é determinístico dado o codificador inicial.

\subsection{O laço}

A cada iteração, o modelo prediz, o usuário simulado escolhe o clique, o marcador é acrescentado, o codificador é reestimado com a plasticidade configurada, e a \ac{IoU} é medida novamente. O laço termina quando a \ac{IoU} ultrapassa o alvo ou quando o teto de cliques é atingido.

O laço possui dois pontos em que os componentes se encaixam. O \ac{AL} pode reordenar os candidatos a clique, em lugar de sempre tomar o maior erro, e o \ac{CL} controla como o clique é incorporado, pelo parâmetro $\alpha$ da Equação \ref{eq_retencao}.

\section{A plataforma interativa}

Os experimentos descritos neste capítulo são executados em lote, com usuário simulado. Para observar o comportamento do sistema com uma pessoa real no laço, foi desenvolvida uma aplicação que roda localmente e integra os três componentes em um fluxo único.

\subsection{O ciclo}

A aplicação opera em rodadas. A cada rodada, o estado consiste no conjunto de marcadores acumulados, no codificador estimado a partir deles, e na curva de qualidade registrada até ali.

O \ac{AL} atua na etapa de sugestão: a aplicação calcula o mapa de probabilidade, pontua as regiões candidatas, e apresenta ao especialista a de maior incerteza. O especialista decide se aceita a sugestão ou se marca outro ponto, e informa se a região é objeto ou fundo. O \ac{IL} atua na incorporação: o toque é materializado como disco de raio 5, acrescentado ao conjunto de marcadores, e o codificador é reestimado. O \ac{CL} atua no controle da reestimação, pelo parâmetro de plasticidade da Equação \ref{eq_retencao}, que pode ser alterado \emph{entre} rodadas.

Essa última possibilidade é deliberada. Alterar a plasticidade no meio de uma sessão permite comparar os dois comportamentos sobre a mesma imagem e o mesmo histórico de cliques, o que torna a diferença visível sem exigir duas sessões separadas.

\subsection{O que a tela registra}

Além da máscara corrente, a aplicação mantém a curva de qualidade por clique e três indicadores: a qualidade inicial, a melhor qualidade alcançada, e o número de cliques que \emph{pioraram} o resultado. O último indicador é o motivo de a tela existir. Com plasticidade máxima, que corresponde ao comportamento original, ele cresce ao longo da sessão, e o especialista observa diretamente o fenômeno descrito na Seção \ref{sec_cl_flim}: anotar mais não implica melhorar.

\subsection{Limite de uso dos números da aplicação}

A qualidade exibida na tela é calculada contra a anotação de referência do conjunto de dados, que existe porque se trata de ambiente de demonstração. Em uso real não haveria essa referência, e o especialista julgaria pela própria máscara.

Os valores apresentados pela aplicação não são usados em nenhuma afirmação desta dissertação. Eles vêm de execução única, sem repetição e sem controle de partição, e servem para inspeção qualitativa. Todas as medições relatadas no Capítulo \ref{cap_4} vêm das campanhas em lote, com semente, partição e registro.

\section{Protocolo experimental}

\subsection{Partições}

Cada conjunto de dados é dividido em três partes disjuntas: um conjunto de onde saem as imagens a anotar, um conjunto de validação usado nas medições de diagnóstico, e um conjunto de teste. A divisão é feita uma vez, com semente fixa, e mantida em todas as campanhas, de modo que resultados de campanhas diferentes sejam comparáveis.

O conjunto de teste não é usado em nenhuma decisão de método. Ele existe para ser tocado uma única vez, na configuração final.

\subsection{Repetições e unidade de análise}

Cada condição é repetida com dez sementes. Uma semente determina quais imagens são sorteadas para anotação e quais traços são gerados sobre elas, mantendo fixos a partição e o conjunto de validação.

A unidade independente de análise é a semente. Quando a mesma condição é avaliada com mais de um decodificador, os decodificadores de uma mesma semente compartilham o codificador e não constituem observações independentes; eles são colapsados por média antes de qualquer teste. Tratá-los como independentes multiplicaria artificialmente o tamanho amostral.

A consequência dessa escolha merece registro explícito: os intervalos de confiança descrevem a variação sob sorteio das imagens de treino \emph{naquele} conjunto de validação, e não generalizam para o conjunto de dados como um todo.

\subsection{Testes estatísticos}

As comparações são pareadas. Dentro de cada célula, os braços compartilham a partição, as imagens e os traços, de modo que a diferença entre eles isola o tratamento. Usa-se o teste $t$ pareado bicaudal, com nível de significância de $0{,}05$.

Quando uma campanha produz múltiplas comparações, aplica-se a correção de Benjamini-Hochberg sobre a família de testes, e o número de comparações que sobrevivem é reportado junto com os valores não corrigidos.

Resultados sobre proporções, como a taxa de sucesso do laço interativo, usam o teste exato de McNemar sobre os pares discordantes, já que pares concordantes não carregam informação sobre a diferença entre os braços.

\subsection{Pré-registro}

Cada campanha é precedida de um documento que declara, antes da execução: a hipótese, o desfecho primário, os desfechos secundários, o critério que falsearia a hipótese, o número de repetições, e o que o experimento não autorizará afirmar. O documento é versionado junto com o código.

Essa disciplina tem motivo prático. Em um espaço com muitos critérios, conjuntos de dados e orçamentos, encontrar uma combinação aparentemente favorável é quase certo mesmo quando nenhum efeito existe. Declarar antes o que contaria como evidência é o que distingue um achado de um acidente. Quatro resultados aparentemente favoráveis foram descartados por esse procedimento ao longo deste trabalho, conforme relatado no Capítulo \ref{cap_4}.

Análises não previstas no pré-registro são permitidas e reportadas, mas identificadas como exploratórias, e qualquer afirmação derivada delas exige confirmação em repetições inéditas.

\subsection{Procedência dos resultados}

Toda execução avaliada é gravada em um registro único, com uma linha por execução, contendo a configuração, a semente, as imagens usadas, as métricas obtidas e o identificador do commit que produziu o resultado. O identificador de cada execução é derivado dos campos que a definem, de modo que reexecutar a mesma configuração produz o mesmo identificador e a duplicata se torna detectável.

As tabelas desta dissertação são geradas a partir desse registro por um programa, sem etapa manual. Cada tabela gerada acompanha um arquivo que lista os identificadores das execuções que produziram cada linha, o que torna a pergunta sobre a origem de um número respondível de forma direta.

\subsection{Reprodutibilidade numérica}

A estimação dos filtros envolve agrupamento por $k$-médias, cuja implementação recorre a bibliotecas de álgebra linear que distribuem o cálculo em várias linhas de execução. A ordem em que as somas parciais são reduzidas varia entre execuções, o que altera o resultado no último dígito representável.

No regime em que o agrupamento não ocorre, descrito pela Equação \ref{eq_regimes}, essa variação não tem efeito. No regime em que ocorre, ela é suficiente para que o agrupamento convirja para uma partição diferente, e o efeito se amplifica entre camadas, já que cada camada recebe como entrada as ativações da anterior.

Todas as campanhas são executadas com as bibliotecas limitadas a uma única linha de execução, condição sob a qual duas execuções idênticas produzem resultados idênticos. A verificação desse comportamento é relatada no Capítulo \ref{cap_4}.

\section{Materiais}

\subsection{Conjuntos de dados}

Três conjuntos são usados, escolhidos por apresentarem propriedades distintas quanto ao tamanho relativo do objeto, que a investigação mostrou ser determinante.

O conjunto de \textbf{parasitas} contém imagens de microscopia de ovos de \textit{Schistosoma mansoni}, em cores, com resolução de $400 \times 400$ pixels. O objeto ocupa cerca de 3\% da imagem. Estão disponíveis marcadores desenhados por dois especialistas, usados nas comparações que envolvem anotação real.

O conjunto \textbf{BraTS} \cite{menze2014multimodal} contém cortes de ressonância magnética cerebral, em tons de cinza, com resolução de $240 \times 240$ pixels.

O conjunto de \textbf{conjuntivite} contém fotografias oculares em cores e alta resolução, nas quais o objeto ocupa fração bem maior da imagem, com área mediana de aproximadamente 89 mil pixels.

\subsection{Levantamento de domínios candidatos}

A escolha dos três conjuntos acima partiu de um levantamento mais amplo de domínios em que a segmentação interativa reduziria esforço de anotação de forma significativa. O critério de seleção considerou a diversidade de desafios que cada domínio impõe, e não apenas a disponibilidade dos dados.

Em gastroenterologia e dermatologia, o Kvasir-SEG \cite{jha2020kvasir} apresenta bordas irregulares de pólipos e artefatos de iluminação, enquanto o ISIC 2018 \cite{codella2018skin} caracteriza-se por alta variabilidade de cor e textura da pele, com interferência de sombras, presença de pelos e heterogeneidade entre pacientes.

Em radiologia e oncologia, o BraTS \cite{menze2014multimodal} e o LiTS \cite{bilic2019liver} introduzem modalidades volumétricas processadas em fatias bidimensionais, com canais monocromáticos e contraste baixo entre tecidos.

Em ultrassonografia e oftalmologia, o BUSI \cite{al2020dataset} impõe ruído característico da modalidade, e o REFUGE2 \cite{fang2022refuge2} exige definição anatômica precisa de estruturas oculares.

Dos domínios levantados, três foram efetivamente usados nos experimentos. A razão é que a investigação revelou, ao longo do trabalho, que a propriedade determinante para o comportamento dos métodos avaliados não é a modalidade de aquisição, e sim a \emph{fração da imagem ocupada pelo objeto}, que controla a razão entre candidatos a filtro e canais disponíveis descrita na Equação \ref{eq_regimes}. Os três conjuntos escolhidos cobrem essa propriedade em faixas bem separadas, de cerca de 3\% a mais de 20\%, o que se mostrou mais informativo do que acrescentar domínios com fração semelhante.

\subsection{Modelos de comparação}

Além do \ac{FLIM} com decodificador adaptativo, foram treinados cinco modelos de referência para contextualizar os valores obtidos: SAMNet, MSCNet, MEANet, UNet e uma variante UNet com codificador \ac{FLIM}. Esses modelos são treinados por retropropagação e com conjuntos anotados completos, de modo que não constituem comparação direta quanto ao esforço de anotação; servem para situar a ordem de grandeza da qualidade atingível em cada domínio.

\subsection{Ambiente}

Os experimentos foram executados em ambiente nativo, com a biblioteca de agrupamento do \textit{scikit-learn}. Uma implementação alternativa de $k$-médias, mais rápida, está disponível na biblioteca original mas não foi utilizada, porque produz centróides diferentes e portanto filtros diferentes, o que tornaria os resultados incomparáveis entre si.

A avaliação não emprega o delineamento por árvores dinâmicas descrito no trabalho de referência, por indisponibilidade do executável na plataforma utilizada. Em seu lugar, usa-se limiarização de Otsu seguida de filtro de área, conforme descrito na Seção \ref{sec_decoder}. Os valores reproduzidos, portanto, não são diretamente comparáveis aos publicados, e essa limitação acompanha toda comparação com o trabalho original.
